\section{Introduction}

Prison overcrowding is a global phenomenon. In 2021, 118 countries reported prison populations exceeding official capacity, and in 11 of these, occupancy rates were more than double the maximum limit. Overcrowding is often attributed to shifts in government policy, institutional omissions, or chronic underinvestment in correctional infrastructure \citep{GlobalPrisonTrends}. Yet, once capacity is exceeded, there is limited causal evidence on how it shapes the behavior of other actors in the criminal justice system—particularly its effects on policing strategies, such as crackdowns, and on judicial sentencing practices.

Understanding these dynamics requires looking beyond formal legal rules. The criminal justice system is often described as the straightforward application of law to individual cases. In practice, however, judges, prosecutors, defense attorneys, and police operate under institutional constraints that shape how they exercise discretion \citep{Ouss-Incentives, Ater-Police-Org-JPE}. Many of these constraints stem from limited resources—whether from rising caseloads or a lack of prison capacity \citep{Schnepel, Yang-Resource-AEJ}. When prisons are overcrowded or dockets overloaded, the incentives and costs facing these actors shift, leading to decisions that can substantially alter sentencing and incarceration rates. In turn, these pressures can produce outcomes that are both ineffective and unjust.

Mexico offers a valuable setting to study this problem. The federal and state prison systems have long faced chronic overcrowding, a problem amplified by the War on Drugs initiated in 2000 under the Partido de Acción Nacional (PAN) government. To meet rising demand, the federal government launched an ambitious expansion, constructing a series of new federal prisons (CEFERESOS) between 2000 and 2011, and in 2012 turning to private contractors under long-term concessions. This paper examines whether the construction of federal prisons during 2000–2012 —representing plausible shocks to local capacity and case-loads— altered state-level judicial sentencing behavior. Specifically, it investigates whether changes originated from shifts in capacity constraints, resource constraints, changes in policing or the heightened salience of incarceration.

To identify the effect of prison construction on state-level sentencing decisions, I rely on variation in incarceration capacity that is plausibly exogenous to underlying trends in sentencing or policing. Specifically, I exploit the staggered construction of new federal prisons across Mexico, which represented shocks to local prison capacity plausibly unrelated to state-level incarceration dynamics. This identification strategy rests on two institutional features: 1) a law requiring judges to assign sentenced individuals to the prison closest to their residence and 2) the fact that state prisons also housed federal level prisoners.\footnote{Article 49 of the \textit{Ley Nacional de Ejecución Penal} (National Law on the Execution of Criminal Sentences) reads: '\textit{... Sentenced individuals may serve their sentence in the prison facilities closest to their place of residence...} \citep{LNEP2024}}

Specifically, I examine the causal effect of prison construction on sentencing outcomes by implementing a staggered difference-in-differences approach, defining as treated those municipalities where the nearest federal prison is between a defined radius (300KM). To consider treatment effect heterogeneity and the presence of negative weights \citep{Chaise-ReviewDID}, I implement \citep{Callaway}, as well as their diagnostic tests to provide evidence on parallel trends (PT) identifying assumptions.

For sentencing information, I use the Criminal Court Proceedings dataset from INEGI, which records cases from criminal and mixed pre-trial courts under state jurisdiction—and in Mexico City, under federal jurisdiction—for 1997–2012. This source includes two complementary datasets. The pre-trial dataset covers all individuals formally processed, capturing early judicial rulings such as pretrial detention or release, which reflect both prosecutorial choices and judicial discretion. The criminal courts dataset, in contrast, documents individuals who received a final decision, including acquittals, prison sentences, or fines. Both datasets provide rich individual-level detail, such as demographics, crime characteristics, sentencing outcomes, and sentencing length. For analysis, I aggregate these records to the municipality level, focusing on variables most relevant to policing and judicial outcomes, such as the number of individuals processed, sent to pretrial detention, released, or sentenced to prison.

To measure prison capacity, I constructed a novel bimonthly panel dataset of all federal and state prisons in Mexico from 2000 to 2014 using over 10,000 OCR-processed pages from the Ministry of Prevention and Social Reintegration. This dataset records openings, closures, capacity, and inmate populations, allowing me to link judicial data to the nearest operational prison for each municipality over time. The treatment variable is a time-varying indicator equal to one in all months after the opening of the closest federal prison within a set distance. 

I show that the construction of a federal prison near a municipality increases both number of individuals processed in pre-trial courts and the number of individuals sentenced. Average Treatment on the Treated (ATT) estimates, however, conceal important timing heterogeneity. Event study results reveal no evidence of differential pre-trends but indicate an inverted U-shaped dynamic: sentencing and processing rises in the short term significantly after prison construction, before gradually declining over the long run (within ten years). Overall, these initial findings provide conclusive evidence that prison construction indeed affects case-load dynamics in the criminal justice system, but does not yet allow me to conclude that it causally affects judge sentencing behavior.

In this context, several mechanisms may influence sentencing decisions. First, if the government expands policing alongside new prison capacity, higher sentencing levels may simply reflect a larger caseload rather than shifts in judicial behavior. In principle, a larger caseload should not alter the rate at which judges impose different types of sentences. However, if increased caseloads create resource constraints, the incentives of judges, prosecutors, and defense attorneys may shift. When bottlenecks arise, all actors have stronger incentives to negotiate plea deals that keep individuals out of prison or reduce sentence lengths, in order to avoid trial \citep{Yang-Resource-AEJ}. This would lower the share of individuals sent to prison while increasing the use of fines.

Second, the effect may operate in the opposite direction through capacity constraints. During periods of overcrowding, judges may have imposed more lenient sentences due to limited space. Once new facilities become operational, sentencing could revert to prior levels, producing an apparent increase. This mechanism aligns with \cite{Ouss-Incentives}, who finds that shifting prison costs from the state to county governments reduced incarceration, and with \cite{Schnepel}, who shows that overcrowding in Texas prisons, combined with a failed ballot initiative to expand capacity, led to more diversion away from incarceration.

Finally, a salience effect may also play a role. Judges, under executive pressure or political incentives, may increase incarceration to justify or fill newly built facilities already financed by taxpayers, a mechanism similar to that proposed by \cite{Dippel}.

I provide additional evidence to disentangle these effects. First, I show that the number of individuals found guilty and sentenced to prison rises significantly after prison construction, while the number sentenced to pay a fine does not. Event-study estimates reveal a positive post-treatment trend for prison sentences and a negative one for monetary penalties. These contrasting results suggest a change in judicial behavior rather than a simple spillover from larger caseloads. By contrast, pretrial detention and releases due to insufficient evidence closely mirror overall caseload dynamics, consistent with the view that rising caseloads are also part of the story. However, there is no evidence that this pattern spilled to criminal courts.

To further examine judicial behavior, I analyze additional outcomes. I find no treatment effects on sentence length (the intensive margin), although event studies indicate a positive but statistically insignificant post-treatment trend. Moreover, both the share of guilty individuals sentenced to prison and the share placed in pretrial detention increase following the opening of a nearby federal prison. While imputing zeros for municipalities without cases introduces some bias, these patterns are more consistent with shifts in judicial decision-making than with caseload growth alone. Finally, I show that prison construction significantly reduces overcrowding at the state level. Event-study results reveal no differential pre-trends but show a long-lasting decline in overcrowding after construction, consistent with a structural increase in capacity. Within six years of a federal prison’s opening, state prisons that housed federal inmates and are on the vicinity of a federal prison experience substantially less overcrowding.

I run several robustness checks to strengthen the claim of my results. First, given the choice of a threshold to define treatment, I provide evidence that the results are robust to different threshold near the 300KM one (-+50KM, with 1KM increases/decreases) as well as defining treatment at the 400KM threshold and 200KM. I also provide evidence that most of the results are robust to other heterogeneity treatment effect (HTE) methods, such as \citep{Gardner, SunAb, Borusyak-DIDImp, Liu-FECT}. Moreover, given that for the levels variables I use a log-transformation of the outcome, and taking into account the potential bias this might introduce in settings with an extensive and intensive margins, I show that the treatment does not affect the extensive margin of the dependent variables, and results are robust to the \cite{ChenRoth2024} transformation that calibrates explicitly the extensive margin.

Finally, the discussion section provides future steps to take to continue the research as well as to provide more robustness to the results. It discusses the necessity to implement a continuous DiD approach, that leverages the variation in distance to the closest federal prison center as the treatment variable using a method that is not biased.

This paper contributes to several strands of the literature. First, it adds to research on the influence of extralegal factors on judicial decision-making, a field largely focused on the U.S. \citep{Posner-Judges, Ouss-Incentives, Schnepel, Gordon-Judges-Election, Kling-Random, Aizer-Juvenile, Berdejo-PoliticalCycles, Dippel, Galinato-2020}. A strong connection to my argument is that of \cite{Ouss-Incentives}, who shows that misaligned incentives between county and state officials alter sentencing behavior.  Also, \cite{Dippel} demonstrates that the construction of private prisons raises sentencing rates through a salience effect. The main contribution of this paper lies in the context: unlike U.S. judges, Mexican judges did not bear fiscal responsibilities nor face electoral incentives. The fact that sentencing behavior still responded to changes in prison capacity makes the findings especially revealing.

I also contribute to the literature on how resource and capacity constraints affect the criminal justice system \citep{Yang-Resource-AEJ, JudgeCaseloads-JPE-Shumway, Robel1990, Huang2011}. Evidence from U.S. courts shows that exogenous increases in caseloads lead prosecutors to dismiss more cases, increase the likelihood of plea deals, and reduce incarceration during judicial vacancies \citep{Yang-Resource-AEJ}. In contrast, this paper shows the opposite pattern: even under exogenous shocks to caseload pressures, judges respond by sending more individuals to prison rather than fewer.

Finally, it contributes to the literature on the driving factors of the global increase in incarceration and sentencing lengths \citep{GlobalPrisonTrends}, again particularly focused in the US \citep{Pfaff2007prison, RaphaelPrisonIncrease, Ouss-Incentives, Mukherjee-Private, Ater-Police-Org-JPE}. Where \cite{Mukherjee-Private} show causally that private incentives inside prison managements lead to a significant increase in sentencing length, through an increase in the probability of getting a conduct violation.  

The remainder of the paper is organized as follows. Section \ref{sec:lit} sets up the theoretical argument and motivates the empirical question. Section \ref{sec:back} outlines the relevant context. Section \ref{sec:data} describes the data. Section \ref{sec:emp} presents the empirical strategy. Section \ref{sec:results} presents the main results, mechanism section, as well as the evidence on the identifying assumptions, robustness checks and discussion. Finally, Section \ref{sec:conclusion} offers concluding remarks.


